\section{Geometría hiperbólica del par angular y recuperación de la acción}
\label{rec:radial}

\subsection{Radión, elipse de acción y coordenatizaciones marcadas}

La sección orientada de fase, acción y memoria denominada \emph{radión} conserva conjuntamente el incremento angular \(\theta=1\), la acción firmada \(\sigma\hbar_{\rm ret}\), la pertenencia a la vuelta \(2\pi_{\HMT}\), la hoja, el retorno nonádico y el cilindro de refinamiento. Su empalme de unidad se expresa en la coordenatización del continuo por \(1=S_E-\Phi_T+\varepsilon_1\). La proyección que retiene solamente \(\theta=1\) produce el radián ordinario; la sección enriquecida conserva además el registro necesario para transportar e invertir sus otras lecturas. En esta sección, \emph{radio} designará una coordenada positiva de forma, distinta de ese objeto enriquecido.

Para explicitar los términos del empalme, la rueda dodecafásica tiene centros \(30m-10\) en su coordenatización de grados. La segunda fase aporta \(50\), y la representación \(A=1000\alpha\) determina
\[
 S_E=\frac{2\pi}{360}(50+1000\alpha).
\]
El sector activo del selector trítico reúne \(H_+=\{1,4,7\}\) y \(H_-=\{2,5,8\}\). Las quince parejas no ordenadas, leídas en las seis posiciones activas, producen \(N_6=6\binom62=90\); su extensión a ocho posiciones produce \(N_8=8\binom62=120\). Las dos hojas del sector activo, normalizadas por la fibra ambiente de \(729\) elementos, fijan
\[
 \rho_{\rm tw}=\frac{2N_6}{729}=\frac{20}{81},\qquad
 \Phi_T=\frac{20}{81}\Delta_4,\qquad
 S_T=1+\Phi_T,\qquad \varepsilon_1=S_T-S_E,
\]
donde \(\Delta_4\) es la coordenada torsional de \eqref{rec:eq:electron-basal}. La resta conserva las secciones que la preceden y prueba directamente
\(S_E-\Phi_T+\varepsilon_1=S_E-\Phi_T+1+\Phi_T-S_E=1\).
El coeficiente procede del censo y de las dos hojas, no del valor de la diferencia final.

Las representaciones angulares de las secciones precedentes satisfacen, en la coordenatización considerada,
\[
 A>0,\qquad |C^*|<A,\qquad
 \xi=\frac{C^*}{A},\qquad \eta_{\rm el}=\operatorname{artanh}\xi.
\]
La notación \(\eta_{\rm el}\) distingue la anisotropía elíptica del coeficiente de retorno \(\eta_{\rm ret}\). Para una sección positiva de acción \(\hbar\), los semiejes son
\begin{equation}
 a_S=\sqrt{2\hbar}\,e^{\eta_{\rm el}/2},\qquad
 b_S=\sqrt{2\hbar}\,e^{-\eta_{\rm el}/2},\qquad a_Sb_S=2\hbar.
 \label{rec:eq:elipse}
\end{equation}
El producto fija la acción por vuelta, mientras el cociente fija la forma. Definimos
\[
 R_\star=\sqrt{b_S/a_S}=e^{-\eta_{\rm el}/2},\qquad
 \tau_{\rm el}=iR_\star^2.
\]

\begin{theorem}[Compatibilidad de las coordenatizaciones radial, de Poincaré y de Klein]
\label{rec:thm:klein}
Sean
\[
 w=\frac{\tau_{\rm el}-i}{\tau_{\rm el}+i},\qquad
 x_K=\frac{2w}{1+w^2}.
\]
En el diámetro real del disco de Poincaré y su imagen en el de Klein,
\begin{align}
 w&=\frac{R_\star^2-1}{R_\star^2+1}
      =-\tanh\frac{\eta_{\rm el}}2,\label{rec:eq:poincare}\\
 x_K&=\frac{R_\star^4-1}{R_\star^4+1}
      =-\tanh\eta_{\rm el}=-\frac{C^*}{A}.\label{rec:eq:klein}
\end{align}
Todas estas transformaciones son invertibles en sus dominios; en particular,
\(R_\star^2=(1+w)/(1-w)\).
\end{theorem}
\begin{proof}
Se sustituye \(\tau=iR_\star^2\) en la transformación fraccionaria. El cociente simplificado tiene denominador \(R_\star^2+1>0\) y \(-1<w<1\). Al sustituir este cociente en \(2w/(1+w^2)\) y multiplicar numerador y denominador por \((R_\star^2+1)^2\), se obtiene \((R_\star^4-1)/(R_\star^4+1)\). La igualdad \(R_\star^4=e^{-2\eta_{\rm el}}\) da la tangente hiperbólica. Finalmente, \(\tanh\eta_{\rm el}=\xi\). Los recuperadores se obtienen despejando los cocientes.
\end{proof}

El signo de \eqref{rec:eq:klein} corresponde a la orientación elegida para el diámetro. La orientación opuesta cambia ambos signos. Es una coordenada hiperbólica de la coordenatización rectangular marcada: conserva la marca y no se confunde con una forma modular ni con una coordenada del cociente modular desprovisto de orientación.

El intercambio de semiejes actúa por
\begin{equation}
 (R_\star,\eta_{\rm el},\tau_{\rm el},w,x_K)
 \longmapsto(R_\star^{-1},-\eta_{\rm el},-1/\tau_{\rm el},-w,-x_K).
 \label{rec:eq:reciproco-radial}
\end{equation}
Sobre este diámetro, las métricas hiperbólicas de curvatura \(-1\) satisfacen
\begin{equation}
 \frac{|d\tau|^2}{(\Im\tau)^2}
 =\frac{4\,dw^2}{(1-w^2)^2}
 =\frac{dx_K^2}{(1-x_K^2)^2}
 =d\eta_{\rm el}^2=4\frac{dR_\star^2}{R_\star^2}.
 \label{rec:eq:metricas}
\end{equation}
Estas igualdades se verifican derivando \eqref{rec:eq:poincare}--\eqref{rec:eq:klein}. En particular, \(d(0,w)=|\eta_{\rm el}|\) y \(d(w,-w)=2|\eta_{\rm el}|\). La transformación \(w\mapsto-w\) es una rotación de \(\pi\) en todo el disco y la simetría respecto del centro en el diámetro generado.

La normalización \((Q/a_S,P/b_S)\) del interior de la elipse produce otra realización del disco de Klein, con coordenadas de posición. Su métrica es
\[
 ds_K^2=\frac{(1-r^2)(du^2+dv^2)+(u\,du+v\,dv)^2}{(1-r^2)^2},
 \qquad r^2=u^2+v^2<1.
\]
Esta coordenatización física y la coordenatización de forma marcada conservan dominios diferentes. En la rama focal \(\eta_{\rm el}\ge0\), la excentricidad \(e_f=\sqrt{1-e^{-2\eta_{\rm el}}}\) da profundidad \(u_f=\operatorname{artanh}e_f\), de donde \(\eta_{\rm el}=\log\cosh u_f\). La profundidad focal no se identifica con la anisotropía.

\begin{figure}[htbp]
\centering
\begin{tikzpicture}[>=Latex,scale=1.0,every node/.style={font=\small}]
\draw[ink,thick] (0,0) circle (1.55);
\draw[->] (-1.8,0)--(1.85,0);
\fill[accent] (-.75,0) circle (2pt) node[below=4pt] {\(w\)};
\fill[accent] (.75,0) circle (2pt) node[below=4pt] {\(-w\)};
\draw[<->,accent,thick] (-.75,.30)--(.75,.30);
\node at (0,-2.0) {Coordenatización de forma: \(w\mapsto-w\)};
\begin{scope}[xshift=7cm]
\draw[ink,thick] (0,0) circle (1.15);
\draw[->] (-1.6,0)--(2.05,0);
\draw[->] (0,-1.4)--(0,1.6);
\draw[accent,thick] (0,0)--(30:.65);
\draw[accent,thick,dashed] (30:.65)--(30:2.03);
\fill[accent] (30:.65) circle (2pt) node[above=3pt] {\(z\)};
\fill[accent] (30:2.03) circle (2pt) node[above=3pt] {\(1/\overline z\)};
\node at (0,-2.0) {Coordenatización espectral: inversión radial};
\end{scope}
\end{tikzpicture}
\caption{Dos reciprocidades tipadas. La primera permanece dentro del disco y transporta la orientación de la forma; la segunda intercambia el interior y el exterior del círculo unidad. Las posiciones dibujadas son esquemáticas.}
\label{fig:rec-dos-cartas}
\end{figure}

\subsection{Sustitución completa de las representaciones angulares}
\label{rec:sustitucion}

La clausura dodecafásica ya construida expresa la misma coordenada \(\alpha\) mediante
\[
 \alpha=\pi_{\HMT}+e_{\HMT}-\varphi_{\HMT}-4-\kappa_{\rm ag},
 \qquad \kappa_{\rm ag}=\frac{K_{\rm int}}{1000^{12}-1}.
\]
Se mantienen la base, el origen de fase y la frontera de acarreo compatibles. La recurrencia analítica completa del artículo representa esa coordenada correlacionada; su representación no se sustituye por un jet finito. Abreviando los subíndices HMT, los lectores de acción de \eqref{iv:action:h5}--\eqref{iv:action:eta-ret} son
\begin{align}
 H_5&=\frac12\sqrt{\frac{1000\alpha}{\varphi}}-\alpha
       +\frac9{16}\alpha^2-\frac59\alpha^3
       +\frac7{48}\alpha^4-\frac1{54}\alpha^5,\nonumber\\
 D_A&=e^{-100\pi\alpha/9},\qquad
 \mathcal C_\pi=169\alpha^6+\frac{D_A\alpha^7}{1-D_A\alpha},\nonumber\\
 \eta_{\rm ret}&=H_5-\frac{90}{\pi}\mathcal C_\pi,
 \qquad A=1000\alpha,\qquad C^*=2(\eta_{\rm ret}+\alpha).
 \label{rec:eq:lectores-completos}
\end{align}
La continuación racional conserva todos los términos de la recurrencia regional. Su convergencia se obtiene sin truncación: para \(\alpha>0\),
\[
 0<\alpha e^{-100\pi\alpha/9}\le\frac9{100\pi e}<1,
\]
pues \(xe^{-cx}\) alcanza su máximo \(1/(ce)\) en \(x=1/c\).

\begin{proposition}[Cancelación estructural en el contraste angular]
\label{rec:prop:sustitucion}
El cociente generado \(\xi=C^*/A\) satisface exactamente
\begin{align}
 \xi={}&\frac1{\sqrt{1000\varphi\alpha}}
 +\frac{9\alpha}{8000}-\frac{\alpha^2}{900}
 +\frac{7\alpha^3}{24000}-\frac{\alpha^4}{27000}\nonumber\\
 &-\frac9{50\pi}
 \left(169\alpha^5+\frac{D_A\alpha^6}{1-D_A\alpha}\right).
 \label{rec:eq:xi-completa}
\end{align}
\end{proposition}
\begin{proof}
Se parte de \(\xi=(\eta_{\rm ret}+\alpha)/(500\alpha)\). El término \(-\alpha\) de \(H_5\) se cancela con \(+\alpha\) del contraángulo. Dividir cada término restante por \(500\alpha\) produce la expresión: la raíz da \(1/\sqrt{1000\varphi\alpha}\), y \(90/500=9/50\). La continuación regional conserva su denominador completo.
\end{proof}

Por tanto, \(x_K=-\xi\) y \(R_\star=((1-\xi)/(1+\xi))^{1/4}\) son expresiones explícitas de las representaciones regionales y del registro. La sustitución de la clausura de \(\alpha\) completa la dependencia de \(\pi,\varphi,e,\kappa_{\rm ag}\), sin promover reconocimientos metrológicos a generadores.

\subsection{Recuperación de Planck con orientación conservada}

\begin{theorem}[Recuperador radial de la sección de retorno]
\label{rec:thm:accion}
En la coordenatización positiva, si \(q=R_\star^4\), entonces
\begin{align}
 \eta_{\rm ret}&=\alpha\frac{499-501q}{1+q},\label{rec:eq:eta-radio}\\
 \hbar_{\rm ret}&=\alpha\frac{499-501q}{1+q}\,10^{-34}\mathcal U_S,
 \qquad\hbar_{\rm pre}=H_5\,10^{-34}\mathcal U_S,\label{rec:eq:planck-radio}\\
 R_{\rm act}&=\frac{\hbar_{\rm pre}}{\hbar_{\rm ret}}
 =\frac{H_5(1+q)}{\alpha(499-501q)}.\label{rec:eq:ratio-radio}
\end{align}
La rama \(\eta_{\rm ret}>0\) equivale a \(q<499/501\), dentro de \(q>0\).
\end{theorem}
\begin{proof}
Las identidades angulares dan \(\eta_{\rm ret}=\alpha(500\xi-1)\). Se sustituye \(\xi=(1-q)/(1+q)\); el numerador es \(499-501q\). El denominador y \(\alpha\) son positivos. La realización dimensional y el cociente son los de las secciones de acción previamente construidas. La década \(-34\) y la unidad \(\mathcal U_S\) conservan su derivación determinantal anterior.
\end{proof}

Esta fórmula explicita la sección de Planck dentro del contraángulo. En términos de las magnitudes de acción, la relación angular es
\[
 C^*=2\left(\frac{\hbar_{\rm ret}}{10^{-34}\mathcal U_S}+\alpha\right).
\]
La división por la base dimensional extrae el coeficiente adimensional; no suma una magnitud de acción a una constante sin unidades.

Para la orientación angular \(\varepsilon\in\{-1,+1\}\), definimos
\[
 C^*_{\varepsilon}=2\varepsilon(\eta_{\rm ret}+\alpha),\quad
 \xi_\varepsilon=C^*_{\varepsilon}/A,\quad
 q_\varepsilon=\frac{1-\xi_\varepsilon}{1+\xi_\varepsilon}>0.
\]
El recuperador sobre ambas coordenatizaciones es
\begin{equation}
 \eta_{\rm ret}=\alpha\left(500\varepsilon
        \frac{1-q_\varepsilon}{1+q_\varepsilon}-1\right).
 \label{rec:eq:eta-orientada}
\end{equation}

\begin{corollary}[Invariancia del recuperador conjugado]
La transformación \((\varepsilon,q)\mapsto(-\varepsilon,q^{-1})\) conserva \(\eta_{\rm ret}\), \(\hbar_{\rm ret}\) y \(R_{\rm act}\).
\end{corollary}
\begin{proof}
La razón \((1-q)/(1+q)\) cambia de signo al invertir \(q\). El producto por \(\varepsilon\) permanece invariante. Se aplica \eqref{rec:eq:eta-orientada} y después las representaciones de acción.
\end{proof}

Si se descarta la orientación y se mantiene indebidamente \(\varepsilon=+1\), la misma expresión produce \(\eta_{\rm ret}(q^{-1})=-\eta_{\rm ret}(q)-2\alpha\). El registro distingue, por tanto, dos operaciones algebraicamente distintas. La orientación angular \(\varepsilon\) tampoco se identifica sin más con el signo simpléctico \(\sigma\): el intercambio de ejes es antisimpléctico, mientras el giro que los intercambia preserva la forma de área.

\subsection{Entrelazamiento de fase y anisotropía}

Sean
\[
 S_\eta=\operatorname{diag}(e^{\eta/2},e^{-\eta/2}),\quad
 J_0=\begin{pmatrix}0&-1\\1&0\end{pmatrix},\quad
 J_\eta=S_\eta J_0S_\eta^{-1}
       =\begin{pmatrix}0&-e^\eta\\e^{-\eta}&0\end{pmatrix}.
\]
Se tienen \(\det S_\eta=1\), \(J_\eta^2=-I\) y
\(C_\eta(\theta)=e^{\theta J_\eta}=S_\eta e^{\theta J_0}S_\eta^{-1}\).
La ley \(S_{\eta+\delta}=S_\delta S_\eta\) prueba el entrelazamiento exacto
\begin{equation}
 C_{\eta+\delta}(\theta)S_\delta=S_\delta C_\eta(\theta).
 \label{rec:eq:entrelazamiento}
\end{equation}
La transformación preserva \(dQ\wedge dP\) y lleva la evolución de fase a la nueva elipse. Con la normalización de \eqref{rec:eq:elipse}, la acción barrida es \(\hbar\theta\), o \(\sigma\hbar\theta\) en la hoja orientada. La fase \(\theta\), la anisotropía \(\eta_{\rm el}\) y el parámetro \(t\) del flujo de \(K\) conservan así funciones diferentes, relacionadas por operadores explícitos.

\subsection{Geodésica radial y realización diferencial de la dirección de K}

Cada radio de \eqref{rec:eq:flujo} determina una coordenatización rectangular y una coordenada hiperbólica:
\[
 \tau_i(t)=ir_i(t)^2,\qquad
 w_i(t)=\frac{r_i(t)^2-1}{r_i(t)^2+1}=\tanh(th_i).
\]

\begin{theorem}[Geodésica de radios compensados]
\label{rec:thm:geodesica}
En el producto de doce diámetros hiperbólicos,
\begin{equation}
 d(w(s),w(t))=2|t-s|\,|\log\rho|,\qquad
 \sum_i\operatorname{artanh}w_i(t)=0.
 \label{rec:eq:geodesica}
\end{equation}
La composición de avances es
\[
 w_i(t+s)=\frac{w_i(t)+w_i(s)}{1+w_i(t)w_i(s)}.
\]
\end{theorem}
\begin{proof}
En coordenadas \(a_i=\log r_i\), la métrica producto es \(4\sum_i da_i^2\). La curva \(a(t)=th\) es recta, con velocidad \(2\|h\|=2|\log\rho|\), y \(\sum_i a_i=0\). Cada diámetro es totalmente geodésico. La distancia en el producto es la raíz de la suma de los cuadrados de las distancias de sus factores, que son \(2|t-s||h_i|\). La identidad de suma nula y la composición se obtienen de \(\operatorname{artanh}w_i=th_i\) y de la fórmula de adición de la tangente hiperbólica.
\end{proof}

La razón regional determina la velocidad de esta realización; \(K\) determina su dirección. Si \(\rho=1\), la trayectoria es constante. Para \(\rho\ne1\), el espacio radial de producto uno tiene dimensión once, con tangente \(V_{11}=\mathbf1^\perp\), y la normal a la trayectoria dentro de él es
\[
 V_{10}(K)=V_{11}\cap u^\perp,\qquad
 P_{10}(K)=P_{11}-\frac{uu^{\mathsf T}}{\|u\|^2}.
\]
El proyector de la proposición~\ref{prop:k-reduccion-diez} adquiere, por tanto, una realización diferencial: elimina la escala uniforme y distingue la dirección generada de sus diez variaciones transversales. Su dominio es este espacio radial, no el espacio real de dimensión veinticuatro de todos los discos.

Con \(D_t=\operatorname{diag}(\operatorname{sech}^2(th_i))\), la expresión en coordenadas \(w\) es \(D_tP_{10}D_t^{-1}\). En efecto, \(D_t\) es el diferencial de \(a\mapsto\tanh a\), y transforma la métrica \(4I\) en \(4\operatorname{diag}((1-w_i^2)^{-2})\). La proyección transportada es ortogonal para esa métrica.

\subsection{Levantamiento reticular y levantamiento rectangular}

Los radios admiten dos realizaciones que conviene conservar explícitas. El operador \(G_t=\bigoplus_i r_iI_2\) escala simultáneamente los dos ejes de cada plano \(A_2\). Conserva el covolumen total y conmuta con la isometría ternaria. La homotecia de cada plano conserva su módulo complejo interno.

La coordenatización \(\tau_i=ir_i^2\) tiene otra realización, de área fija:
\[
 T_i=\operatorname{diag}(r_i^{-1},r_i),\qquad
 (\omega_1,\omega_2)=(r_i^{-1},ir_i),\qquad
 \omega_2/\omega_1=ir_i^2.
\]
Es el levantamiento elíptico con \(\eta=-2\log r_i\). Esta convención transforma períodos físicos. La acción fraccionaria lineal de la misma matriz sobre \(i\) daría \(i/r_i^2\); para obtener \(ir_i^2\) bajo esa otra convención se utiliza \(T_i^{-1}\). La distinción de acciones evita una inversión inadvertida.

Ambas realizaciones se reconstruyen desde los mismos radios, pero la conmutación ternaria de la homotecia no se transfiere a la deformación anisótropa. El estado y sus lectores permiten comparar las dos geometrías sin identificar sus operadores.
